\documentclass[10pt,a4paper]{article}
\usepackage[margin=1.75cm]{geometry}
\usepackage{graphicx,booktabs,titlesec,hyperref,caption,siunitx,xcolor,amsmath}
\usepackage[T1]{fontenc}
\usepackage[protrusion=true,expansion=false]{microtype}
\hypersetup{colorlinks=true,urlcolor=blue!55!black,linkcolor=black}
\captionsetup{font=small,labelfont=bf,skip=3pt}
\titlespacing*{\section}{0pt}{1.1ex}{0.35ex}
\titlespacing*{\subsection}{0pt}{0.8ex}{0.25ex}
\titleformat{\section}{\normalfont\normalsize\bfseries}{\thesection}{0.55em}{}
\titleformat{\subsection}{\normalfont\small\bfseries}{\thesubsection}{0.5em}{}
\setlength{\parskip}{0.25ex}
\pagestyle{empty}
\sisetup{group-separator={,},group-minimum-digits=4}

\title{\vspace{-2.6em}\textbf{Supply, Demand and Policy in New York's Ridehail Market:\\
A PySpark MapReduce Analysis of 462 Million Trips}\\[0.2em]
\normalsize MIT 805 Big Data --- Part 2: MapReduce \& Visualisation}
\author{\normalsize Gishon K.\ Gwenzi (\texttt{21731056}) \and
        \normalsize Letsoba Savannah Mabe (\texttt{20582995})}
\date{\vspace{-1em}\small University of Pretoria \quad|\quad 8 October 2026}

\begin{document}
\maketitle
\vspace{-2.9em}
\begin{center}\footnotesize
\textbf{Repository:} \url{https://github.com/caie007/savannahgishon}
\end{center}
\vspace{-0.5em}

\section{Processing objective and analytical question}

New York's high-volume for-hire vehicle (HVFHV) market moved \num{461855376}
trips across 2024 and 2025. Part~1 established, on 2025 data alone, that
passenger wait peaks at 04:00 and troughs at 10:00 --- the inverse of the demand
profile --- implying that service quality is governed by driver supply rather
than road congestion. It also showed that the central business district (CBD)
congestion fee introduced on 5~January~2025 is a flat \$1.50 charge, and
therefore proportionally heavier on low-fare trips.

Part~2 tests both propositions at scale and joins them into one question:

\begin{quote}
\emph{Where, when and for whom does New York's ridehail market fail to match
supply to demand --- and did congestion pricing change that?}
\end{quote}

Three sub-questions follow. \textbf{(i)} Did CBD-bound volume fall after the
charge, relative to trips that never enter the zone, and beyond any pre-existing
trend? \textbf{(ii)} Who bore the cost of a flat fee? \textbf{(iii)} Which zones
and hours are persistently under-served, and does supply density explain it?

Answering these requires a full pass over 24 months of trip records --- roughly
\SI{107}{\giga\byte} uncompressed --- which is not tractable on a single-threaded
framework. The analysis is therefore implemented as MapReduce-style distributed
processing in PySpark.

\section{Data and distributed architecture}

The working set is 24 monthly Parquet files (2024-01 to 2025-12),
\SI{11.1}{\giga\byte} compressed. After the Part~1 cleaning rules --- positive
distance and duration, duration $\leq$\,6\,h, non-negative fare, dropoff after
pickup, and both endpoints among the 265 mapped taxi zones --- \num{229391984}
records remain for 2024 and \num{232463392} for 2025.

A single distributed pass reduces these \num{461855376} rows to an aggregate of
\num{73973527} rows keyed by
$(\text{PU zone},\text{DO zone},\text{year},\text{month},\text{dow},\text{hour})$,
storing \emph{sums} rather than means so that any later regrouping re-aggregates
correctly. Every subsequent result in this report is computed from that
aggregate. This pre-aggregation pattern is the practical core of the work: it
converts a problem that does not fit in memory into one that does, and it is the
shared layer from which both authors' threads are derived.

\subsection{Deriving the charged zone set from the data}

The charged area was not hard-coded. In 2025 the \texttt{cbd\_congestion\_fee}
column records which trips were charged, so a map--shuffle--reduce over
destination zone yields the empirical charge rate per zone. The distribution is
strongly bimodal (Appendix Fig.~A1), so a threshold of \num{0.80} is not
delicate: it identifies \textbf{38 CBD zones}, which were then applied backwards
to 2024 to label trips that \emph{would have been} charged. A ring of 32
Manhattan zones immediately outside the charged area was labelled
\emph{boundary} to test for displacement.

\section{MapReduce implementation}

\subsection{Mapping and transformation}
Each record is projected to the columns the analysis needs and enriched with
derived fields: \texttt{year}, \texttt{hour}, a boolean \texttt{cbd\_trip} from
the derived zone set, and
$\texttt{wait\_s} = \texttt{pickup\_datetime} - \texttt{request\_datetime}$, with
a \texttt{wait\_valid} guard excluding the \SI{1.15}{\percent} of records whose
pickup precedes the request (a telematics clock-drift artefact identified in
Part~1). All of this is lazy: the physical plan in Appendix~A shows Catalyst
collapsing six chained projections into two whole-stage code-generated scans.

\subsection{Shuffle and grouping}
Aggregation by zone and hour requires rows sharing a key to meet on one
executor, so Spark repartitions by hash of the grouping key. This is the
\texttt{Exchange hashpartitioning} node in the physical plan and the only point
at which data crosses the network. For the main aggregation the exchange moved
\textbf{\SI{7.2}{\gibi\byte}} --- \SI{6}{\percent} of the \SI{120.5}{\gibi\byte}
read from disk, the reduction coming from column pruning pushed into the Parquet
reader and from partial aggregation applied before the shuffle.

\subsection{Reduction and aggregation}
We implemented the identical computation --- trip count and mean wait by
$(\text{PU zone},\text{hour})$ --- twice, to make the execution models
comparable on like terms (Table~\ref{tab:dual}; see also Fig.~A5).

\begin{table}[h]\centering\small
\vspace{-0.3em}
\caption{Two implementations of one aggregation, \num{1964006} rows, 8 partitions.}
\label{tab:dual}
\vspace{-0.3em}
\begin{tabular}{@{}lrrr@{}}
\toprule
Implementation & Runtime & Shuffle write & Keys produced \\
\midrule
RDD \texttt{map} $\to$ \texttt{reduceByKey} & \SI{12.7}{\second} & \SI{558.3}{\kibi\byte} & \num{6143} \\
DataFrame \texttt{groupBy} $\to$ \texttt{agg} & \SI{1.9}{\second} & \SI{1061.4}{\kibi\byte} & \num{6143} \\
\bottomrule
\end{tabular}
\vspace{-0.3em}
\end{table}

Both produced \num{6143} identical keys, with zero difference in counts and mean
waits agreeing to floating-point precision. The RDD path makes the three phases
explicit: \texttt{map} emits $((\text{zone},\text{hour}),(1,\text{wait}))$,
\texttt{reduceByKey} combines locally on each partition before shuffling --- the
exact function of a Hadoop Combiner --- and a final \texttt{map} converts partial
sums to means. The stage metrics quantify that combining: the \texttt{reduceByKey}
stage read \SI{166.7}{\mebi\byte} and wrote only \SI{558.3}{\kibi\byte} of
shuffle, a \textbf{306-fold reduction} before any data crossed the network.

The decisive observation is that the RDD path shuffled \emph{less} data than the
DataFrame path yet ran \num{6.56} times slower. The gap is therefore not shuffle
volume; it is the process boundary. Python RDD lambdas require every row to be
serialised out of the JVM into a Python worker and back, whereas DataFrame
operations are compiled to JVM bytecode by Tungsten and never leave it. This was
not merely a performance difference: attempting the RDD benchmark over the full
2025 set caused Python worker handshake timeouts, which is the same cost
appearing as a failure rather than a delay.

\section{Distributed execution evidence}

\begin{table}[h]\centering\small
\vspace{-0.3em}
\caption{Spark UI metrics for the full run (application \texttt{MIT805-Part2-HVFHV},
Spark 3.5.1, Java 17.0.20.1, \texttt{local[8]}, \SI{32}{\giga\byte} driver).}
\label{tab:exec}
\vspace{-0.3em}
\begin{tabular}{@{}lr lr@{}}
\toprule
Metric & Value & Metric & Value \\
\midrule
Tasks completed / failed & \num{849} / \textbf{0} & Total input read & \SI{120.5}{\gibi\byte} \\
Cores & 8 & Shuffle read / write & \SI{7.2}{\gibi\byte} / \SI{7.2}{\gibi\byte} \\
Cumulative task time & \SI{44}{\minute} (GC \SI{34}{\second}) & Disk used & \SI{28.2}{\gibi\byte} \\
Jobs / SQL queries & 49 / 15 & Shuffle partitions (AQE) & 96 \\
Stages completed / skipped & 50 / \textbf{30} & Wall clock & \SI{39}{\minute} \\
\bottomrule
\end{tabular}
\vspace{-0.4em}
\end{table}

Three observations from Table~\ref{tab:exec} bear directly on the execution
model. First, \textbf{30 of 80 stages were skipped} because Spark reused shuffle
output already materialised by earlier jobs; a chained MapReduce pipeline has no
equivalent and would recompute from input each time. Second, adaptive query
execution coalesced the configured 256 shuffle partitions down to \textbf{96} at
runtime based on observed partition sizes. Third, the cached intermediate frame
reports \SI{26.1}{\gibi\byte} on disk and \SI{0}{\byte} in memory against
\SI{19}{\gibi\byte} of storage memory: it did not fit, and Spark degraded to disk
rather than failing.

\subsection{Spark's DAG versus Hadoop's Map $\to$ Shuffle $\to$ Reduce}

Hadoop MapReduce executes one fixed Map $\to$ Shuffle $\to$ Reduce cycle per
job and materialises output to HDFS between jobs, replicated three times by
default. Spark instead compiles the whole computation into a DAG of stages whose
boundaries fall only at shuffles, pipelining everything in between in memory.
Four concrete contrasts from this run:

\begin{itemize}\itemsep1pt \parskip0pt
\item \textbf{Intermediate I/O.} The \SI{7.2}{\gibi\byte} that crossed our single
shuffle would, under Hadoop, additionally be written to and re-read from HDFS
between every stage --- roughly \SI{21.6}{\gibi\byte} of extra write traffic at
default replication, none of which Spark incurs.
\item \textbf{Automatic combining.} Catalyst placed a partial
\texttt{HashAggregate} before the exchange and a final one after it. Hadoop
requires a Combiner to be written by hand to achieve the same effect.
\item \textbf{Join strategy.} The 265-row zone lookup was broadcast to every
task, eliminating the shuffle of the large side entirely. Classical MapReduce has
no broadcast join and must shuffle both inputs.
\item \textbf{Lazy, replanned execution.} Nothing ran until an action, and AQE
re-planned partitioning mid-flight. MapReduce is eager and fixed at submission.
\end{itemize}

\section{Results}

\subsection{Thread A --- Congestion pricing did less than the headline suggests}

A naive year-on-year comparison looks dramatic: CBD-bound trips fell
\SI{4.39}{\percent} (\num{79897901} to \num{76390002}) while non-CBD trips rose
\SI{4.40}{\percent}, a difference of \textbf{\num{-8.79} percentage points}. The
CBD share of all trips fell from \SI{34.83}{\percent} to \SI{32.86}{\percent}.

That comparison is misleading, and Fig.~\ref{fig:vol}(b) shows why. The CBD share
was \emph{already} falling through 2024 at \num{-0.103} percentage points per
month, before any charge existed. There is no discontinuity at the policy date:
the share in January 2025 was \SI{+0.21}{} pp \emph{above} December 2024, within
the ordinary month-to-month variation of \num{-0.199} pp observed across 2024.
Extrapolating the 2024 trend predicts a 2025 mean of \SI{33.60}{\percent} against
an actual \SI{32.88}{\percent} --- a residual of only \textbf{\num{-0.72} pp},
and the 2025 slope (\num{-0.125} pp/month) is only marginally steeper than 2024's.

\begin{figure}[h]\centering
\includegraphics[width=0.93\linewidth]{../figures/f1_volume.png}
\vspace{-0.4em}
\caption{Trip volume by charge status. The divergence in (a) is real, but (b)
shows most of it predates the 5 January 2025 charge.}
\label{fig:vol}
\vspace{-0.5em}
\end{figure}

Two further results corroborate a weak volume effect. Trips terminating in the
boundary ring grew \SI{0.79}{\percent} year-on-year against \SI{5.26}{\percent}
for outer zones (Fig.~A2): had riders been avoiding the charge by stopping short
of the boundary, the ring would have gained disproportionately, and it did not.
Meanwhile CBD mean fare rose from \$31.24 to \$33.57 (+\SI{7.5}{\percent}) while
non-CBD fare was flat (\$20.62 to \$20.65). The \$1.50 charge accounts for
roughly \SI{64}{\percent} of that \$2.33 increase, so the cost was passed to
riders rather than suppressing the trip.

\subsection{Who paid}

Conditional on being charged, the flat fee is regressive
(Fig.~\ref{fig:fee}). Against mean borough fares it represents
\SI{7.98}{\percent} of a Bronx fare but only \SI{5.17}{\percent} of a Manhattan
one --- a \num{1.54}$\times$ spread, falling hardest on the lowest-fare riders.
In aggregate the incidence is reversed: Manhattan pickups generated
\SI{83.1}{\percent} of the \$\num{117124706} collected in 2025, because
\SI{76.3}{\percent} of Manhattan trips are charged against \SI{3.7}{\percent} of
Bronx trips. Both statements are true and they answer different questions:
Manhattan funds the scheme, but an outer-borough rider who does cross the cordon
pays half again as much of their fare for the privilege.

\begin{figure}[h]\centering
\includegraphics[width=0.86\linewidth]{../figures/f3_fee.png}
\vspace{-0.4em}
\caption{Flat-fee incidence. Per charged trip the burden is regressive (a);
in aggregate Manhattan funds the scheme (b).}
\label{fig:fee}
\vspace{-0.5em}
\end{figure}

\subsection{Thread B --- Service quality is a function of supply density}

Across \num{5898} zone--hour cells we computed a \emph{mismatch index}: mean wait
in a zone-hour divided by the citywide mean wait for that same hour. Normalising
by hour is essential --- without it the analysis merely rediscovers that 04:00 is
slow everywhere. The index reaches \num{2.35} at its worst, and its 90th
percentile is \num{1.297}.

The correlation between zone trip density and mean wait is \textbf{negative at
every one of the 24 hours} (Fig.~\ref{fig:corr}), from \num{-0.24} at 08:00 to
\num{-0.617} at 20:00, with an overall $r = \num{-0.472}$. Busier zones are
served \emph{faster}, at all times of day. This confirms Part~1's inference
directly: the binding constraint is the density of available drivers, not road
congestion. If congestion dominated, the sign would invert at peak hours, and it
never does.

\begin{figure}[h]\centering
\includegraphics[width=0.60\linewidth]{../figures/f4_corr.png}
\vspace{-0.4em}
\caption{Correlation between zone trip density and mean wait, by hour.}
\label{fig:corr}
\vspace{-0.5em}
\end{figure}

The same mechanism governs shared-ride matching. The match rate ranges from
\SI{70.9}{\percent} at 23:00 to \SI{39.4}{\percent} at 11:00 and tracks request
volume with $r = \num{+0.59}$ (Fig.~A3): thin demand cannot be pooled, so
passengers who accept a discounted shared fare at midday ride alone three times
in five.

Persistence matters more than any single bad hour. Filtering to zone-hours above
the 90th-percentile mismatch and counting distinct bad hours per zone yields
\textbf{74 persistent service deserts}, four of which are above threshold in all
24 hours (Fig.~\ref{fig:des}, Fig.~A4). They are overwhelmingly peripheral ---
23 in Queens, 18 in Staten Island, 18 in Brooklyn --- and concentrated in the
Rockaway peninsula and southern Staten Island. Hammels/Arverne is worst at a mean
index of \num{1.61}, meaning its riders wait \SI{61}{\percent} longer than the
citywide norm for the same hour, every hour of the day.

\begin{figure}[h]\centering
\includegraphics[width=0.66\linewidth]{../figures/f7_deserts.png}
\vspace{-0.4em}
\caption{Zones above the 90th-percentile mismatch index for the greatest number
of hours. Colour denotes borough.}
\label{fig:des}
\vspace{-0.5em}
\end{figure}

\subsection{Where the two threads meet}

The threads converge on one number. Between 2024 and 2025 mean wait in the CBD
\emph{improved} from \SI{272.0}{\second} to \SI{269.2}{\second}
($-\SI{1.0}{\percent}$), while mean wait everywhere else \emph{deteriorated} from
\SI{281.0}{\second} to \SI{294.1}{\second} ($+\SI{4.7}{\percent}$). Congestion
pricing did not materially remove trips from the cordon. What it appears to have
done is redistribute driver supply --- and given that service quality is
determined by supply density at every hour of the day, a marginal gain inside the
cordon was accompanied by a loss four times larger outside it, borne by the same
peripheral zones already identified as service deserts.

\section{Business and societal interpretation}

\textbf{For the TLC and the MTA}, the headline finding is that the fee raised
\$\num{117124706} in its first year with no volume reduction distinguishable from
the pre-existing trend. If the policy objective was revenue, it succeeded; if it
was modal shift away from for-hire vehicles entering the CBD, the evidence here
does not support it. The \num{-0.72} pp residual against the 2024 trend is the
honest upper bound on the volume effect, and even that is not separable from
other 2025 changes.

\textbf{For equity regulators}, the flat structure produces a \num{1.54}$\times$
spread in proportional burden between the lowest- and highest-fare boroughs. A
fee scaled to fare, or a rebate for trips originating in the zones identified in
Fig.~\ref{fig:des}, would raise comparable revenue with materially less
regressivity --- and the dataset makes the trade-off directly computable.

\textbf{For outer-borough riders}, the service-desert result has a sharper
implication than slow service. Part~1 showed these are the zones where ridehail
behaves as transit substitution rather than discretionary travel. They are now
shown to be both the worst-served at every hour and the places where waits
deteriorated most during the policy year. A rider in Hammels/Arverne waits
\SI{61}{\percent} longer than the city norm with weaker rail alternatives than a
rider in Midtown.

\textbf{For operators}, the negative density--wait correlation at every hour is a
positioning signal: the marginal value of an additional driver is highest exactly
where current density is lowest. The 74 zones in Fig.~A4 are a concrete target
list, and the midday shared-match collapse to \SI{39.4}{\percent} identifies when
pooling incentives fail.

\textbf{Limitations.} Three bound these conclusions. \emph{First, identification}
--- 2024 and 2025 differ in many ways besides the fee, and we report association
with an explicit pre-trend adjustment, not causation. \emph{Second,
survivorship} --- the data records only matched, completed trips, so a service
desert cannot be distinguished from suppressed demand; riders who stopped
requesting are invisible, which means our mismatch index understates the problem.
\emph{Third, no entity identifiers} --- without driver or vehicle keys, the
supply redistribution we infer from wait times cannot be confirmed by observing
drivers directly.

\clearpage
\section*{Appendix}

\begin{figure}[h]\centering
\includegraphics[width=0.60\linewidth]{../figures/f0_threshold.png}
\caption*{\small\textbf{Figure A1.} Empirical charge rate by destination zone,
2025. The distribution is strongly bimodal --- 178 zones below 0.2 and 38 above
0.8, with only 29 in between --- so the classification threshold is not delicate.}
\end{figure}

\begin{figure}[h]\centering
\includegraphics[width=0.78\linewidth]{../figures/f2_segments.png}
\caption*{\small\textbf{Figure A2.} Trip volume by zone segment, indexed to each
segment's 2024 mean. The boundary ring grew \SI{0.79}{\percent} against
\SI{5.26}{\percent} for outer zones, which is inconsistent with riders stopping
short of the cordon to avoid the charge.}
\end{figure}

\begin{figure}[h]\centering
\includegraphics[width=0.62\linewidth]{../figures/f6_shared.png}
\caption*{\small\textbf{Figure A3.} Shared-ride match rate against request volume
by hour. Matching succeeds where demand is thick and collapses at midday.}
\end{figure}

\begin{figure}[h]\centering
\includegraphics[width=0.84\linewidth]{../figures/f5_mismatch.png}
\caption*{\small\textbf{Figure A4.} Mismatch index by zone and hour for the 18
worst-served zones. Values above 1.0 indicate waits longer than the citywide mean
for that hour.}
\end{figure}

\begin{figure}[h]\centering
\includegraphics[width=0.70\linewidth]{../figures/f8_rdd_df.png}
\caption*{\small\textbf{Figure A5.} The RDD implementation shuffled less data than
the DataFrame implementation yet ran \num{6.56}$\times$ slower, isolating the
JVM-to-Python process boundary as the cost.}
\end{figure}

\vspace{0.5em}
\noindent\textbf{Physical plan (abridged).} Catalyst collapsed six chained
projections into two whole-stage code-generated scans, pushed the zone predicate
into the Parquet reader as an \texttt{INSET} filter, and placed the exchange as
the single shuffle boundary:

{\footnotesize
\begin{verbatim}
== Physical Plan ==
AdaptiveSparkPlan isFinalPlan=false
+- HashAggregate(keys=[PULocationID#7, _groupingexpression#3145],
                 functions=[count(1), avg(wait_s#343L)])
   +- Exchange hashpartitioning(PULocationID#7, _groupingexpression#3145, 96),
                ENSURE_REQUIREMENTS, [plan_id=1044]
      +- HashAggregate(keys=[...], functions=[partial_count(1), partial_avg(...)])
         +- InMemoryRelation [...], StorageLevel(disk, memory, deserialized, 1 replicas)
            +- *(1) ColumnarToRow
               +- FileScan parquet [...] PushedFilters: [IsNotNull(trip_miles), ...]
\end{verbatim}}
\vspace{-0.6em}
\noindent Full plan: \texttt{output/p2/physical\_plan\_groupby.txt} in the repository.

\end{document}
